\section{Introduction}
\label{sec:intro}

Majority voting over sampled reasoning traces, known as Self-Consistency \citep{wang2023sc}, is one of the most effective and simplest methods for improving large language model (LLM) reasoning accuracy. SC selects the most frequent answer among multiple chain-of-thought samples, and despite its simplicity, it remains competitive across mathematical, logical, and commonsense reasoning tasks \citep{chen2024usc, taubenfeld2025cisc}. A natural question follows: can we extract richer signal from these traces beyond vote counting? Recent work explores this direction by formalizing reasoning steps into logical constraints and using automated solvers for answer selection. Systems such as Logic-LM \citep{pan2023logiclm}, LINC \citep{olausson2023linc}, DyLoS \citep{xu2026dylos}, Aristotle \citep{li2025aristotle}, FoVer \citep{wang2025fover}, and LogicAgent \citep{luo2025logicagent} combine LLMs with symbolic solvers and report gains on logical reasoning benchmarks. These results suggest that formal constraint structure, when properly obtained, can improve upon majority voting.

The practical question for practitioners is not whether symbolic methods can ever help, but when the investment is justified. Adding constraint extraction and solver invocation to the inference pipeline introduces substantial overhead: approximately 48\% more compute in our measurements, plus the engineering complexity of integrating formal solvers. If the gains are inconsistent or limited to specific domains, the additional cost may outweigh the benefit. This question is made more urgent by recent findings that SC exhibits scaling plateaus, where additional sampled traces yield diminishing accuracy returns \citep{loo2025reevaluating}. The plateau motivates the search for better aggregation methods, but it also raises the bar: any proposed improvement must reliably outperform a baseline that is both minimal in complexity and difficult to beat.

Existing neuro-symbolic reasoning systems sidestep a fundamental question by design. Logic-LM \citep{pan2023logiclm} and LINC \citep{olausson2023linc} delegate formalization to external provers. DyLoS \citep{xu2026dylos} routes problems to specialized solver paradigms. FoVer \citep{wang2025fover} trains a dedicated verifier on supervised first-order logic (FOL) data. Aristotle \citep{li2025aristotle} applies logic-complete decomposition with external validation. All of these methods introduce knowledge or capabilities beyond the reasoning model itself. The simplest and most scalable alternative is \emph{same-model self-extraction}: the LLM that generated the reasoning traces also formalizes constraints from those traces, with no additional training data, no external prover, and no second model. This approach requires minimal infrastructure and generalizes to any domain without supervised annotations. Yet no prior work has systematically evaluated whether same-model self-extraction produces constraints useful for answer aggregation, or characterized the conditions under which it succeeds or fails.

We address this gap with Symbolic Invariant Consensus Aggregation (SICA), a pipeline that extracts FOL constraints from sampled reasoning traces and aggregates them through weighted partial MAX-SAT optimization using Z3 \citep{demoura2008z3}. We conduct the first systematic evaluation of same-model self-extraction for constraint-based voting. Our experiments span five models (Qwen2.5-14B, Mistral-7B, LLaMA-3.1-8B, DeepSeek-R1-7B, QwQ-32B), four reasoning domains (FOLIO, ProofWriter, StrategyQA, ARC-Challenge), and more than 1,300 instances. The central finding is negative: SICA provides no statistically significant improvement over SC in any condition, with accuracy gaps between $-3.3$ and $+2.0$ percentage points. We identify the structural cause as the \emph{capability trap}: because the same model both reasons and formalizes, the extracted constraints are semantically coupled to their source answers, and MAX-SAT optimization reduces to simple vote counting. A constraint ablation confirms this directly: MAX-SAT and a count-only baseline agree on 803 of 804 predictions (\S\ref{sec:ablation}). Crucially, an oracle control that replaces self-extracted constraints with gold FOL annotations achieves $+8.82$ percentage points over SC ($p < 0.001$); the aggregation mechanism works when supplied with discriminative constraints (\S\ref{sec:oracle}). The failure is specific to self-extraction, not to symbolic aggregation as a method.

Our contributions are as follows:
\begin{itemize}
    \item \textbf{Systematic negative result.} We report the first controlled evaluation of same-model constraint self-extraction across five models, four domains, and more than 1,300 instances. No condition produces a statistically significant improvement over majority voting (\S\ref{sec:main-results}).
    \item \textbf{Causal isolation via oracle control.} By replacing only the constraint source while holding all pipeline components fixed, we establish that the MAX-SAT aggregation mechanism is sound and that the bottleneck lies entirely in extraction quality. Gold FOL constraints yield $+8.82$\,pp over SC, a recoverable gap of practical significance (\S\ref{sec:oracle}).
    \item \textbf{The $C_0$ diagnostic.} We introduce the multi-candidate rate $C_0$, a zero-cost metric computable before constraint extraction, as a necessary-condition screening criterion. When $C_0 < 5\%$, constraint-based methods degenerate to majority voting by construction (\S\ref{sec:c0}).
    \item \textbf{Constraint ablation.} We show that MAX-SAT, count-only, and random-weight aggregation produce nearly identical predictions, self-extracted constraints carry no discriminative signal beyond vote counts (\S\ref{sec:ablation}).
\end{itemize}
